\subsection{Substitution rung}

\begin{theorem}[Jeandel--Rao marker-domino substitution chain map ($b_3 \mapsto \varphi$)]\label{thm:jeandel-rao-substitution}\lean{Obsmetrologia.AperiodicOrder.Tiling.Witness.jr_substitution_pullback}\leanok\uses{thm:jeandel-rao-witness, thm:jeandel-rao-b3}
The marker-domino step is an explicit integer cellular chain map
$\mathrm{jrComplex}\to\mathrm{markerRingComplex}$. On $1$-cells, the marked
edges map $e_2\mapsto q_1$ and $e_3\mapsto q_5$, while the four unmarked edges
map to $q_0$; on $2$-cells every witness face maps to $-g_3$, forced because
$\partial(-g_3)=q_1-q_5$ is the edge-image of each witness-face boundary. The
induced cochain pullback sends the marker-ring core class to the decorated
witness exactly, $\mathrm{subPull}\,b_3=\varphi$, and carries cocycles to
cocycles. Thus $\varphi$ is the pullback of a nonzero $H^1$ class along a
genuine integer chain map: one marker-domino rung.
\end{theorem}
